% Limitations and Future Work. Four candid groups plus one limitation-and-remedy
% table. The credibility of the bill depends on stating plainly what the record
% does and does not establish.

The credibility of this bill depends on an honest account of what the record does
and does not establish. Four groups of limitations follow; each limitation that
admits a concrete remedy is paired with one in Table~\ref{tab:limits-remedies}.

\subsection{Codegen and Execution Limitations}

Several capabilities could not be run live and were handled by deterministic
fallbacks, and the bill rests on the record as it actually is. The agentic and
on-premise large language model backends were not live and ran on deterministic
offline fallbacks, so intent generation used a frozen reference plan rather than a
live model. Live web search and PDF ingestion were unreachable and used offline
stubs or size estimates. The physics and robotics backends, including the major
simulation engines, were not installed, so behaviors ran on standard-library
references rather than a full physics engine. 

Local execution used a single
Python interpreter, with cross-version coverage delegated to the continuous
integration matrix across Python 3.10, 3.11, and 3.12. In the first study, a
multibyte autochunk step lost six characters on one oversized line, a defect
reported and left in scope rather than repaired. And the first study had no
independent gold-standard validation reference, so one accept case was
self-referential. Two assurance gaps from the first study are therefore on the
record: the absence of an independent reference, and a rich input corpus not yet
wired end to end into the executed pipeline \cite{kawchak2026vvuq01,
kawchak2026vvuq02}.


\subsection{Simulation Caveats}

The Unitree H2-Surgical 1.0 is a clearly labelled hypothetical 2030 platform, and
every supporting number in this bill is a simulation result. The four-entrant
comparison is simulation against simulation, and the human-surgeon baseline round
carries a time-dimension caveat, because it compares a short robot run against a
multi-hour human procedure. A real deployment would require, at minimum, the
robotic-surgery and medical-electrical safety standards, the service-robot
stability standard, Food and Drug Administration Software as a Medical Device
Class III clearance, institutional review board approval, and regulatory
authorization; the bill does not presume any of these \cite{iec8060127, iec60601,
iso13482}.

\subsection{This Full-Paper Step}

This full-paper step has its own limits worth stating. A PDF compiler was not
available in the build container, and the prompt asked that the bill not be
compiled here, so the LaTeX was validated structurally rather than rendered: the
checks confirmed balanced braces and environments, that every citation key is
present in the bibliography, balanced dollar pairs, and single-hyphen and
section-symbol scans, and the bundle was made compile-safe. The rendered PDF and
the final page balancing are left to Overleaf. No images are used anywhere
\cite{lamport1994latex}.

\subsection{Legislative Limitations and Future Work}

This is an independent draft bill, not enacted law, and is not endorsed by CFR,
ICH, FDA, or any sponsor, contract research organization, site, institutional
review board, regulator, or medical society. It reflects one author and needs the
stakeholder input the policy memorandum names. The state-law landscape is moving,
including the Texas committee substitute and the Florida 2026 bills, so the
citations should be refreshed at enactment. For future work, the central forward
claim is that this verification-before-generation standard should extend across
the full breadth of oncology trials, with the remedies in
Table~\ref{tab:limits-remedies} prioritised: supply an independent validation
reference, wire the standards corpus end to end, enable the live backends, repair
the multibyte chunker, and broaden the simulation set toward real-world
validation \cite{kawchak2026bills, kawchak2026prediction, fda2026realtime}.

\begin{table}[ht]
\centering
\begin{tabularx}{\textwidth}{L{7cm} Y}
\toprule
\textbf{Limitation (abbreviated source)} & \textbf{Concrete remedy} \\
\midrule
\multicolumn{2}{l}{\textbf{papers/VVUQ-01/, papers/VVUQ-02/}} \\
LLM backends offline (exec. READMEs) & Enable live backends; re-run with the same seed \\
No ind. reference (VVUQ-01 limits) & Supply gold-standard ref.; retest validation \\
Offline, PDF ingestion (exec. READMEs) & Connect live retrieval; re-ingest the full corpus \\
No physics engine (VVUQ-02 limits) & Run behaviors under a full physics engine \\
Single interpreter locally & Already covered by the CI 3.10 to 3.12 matrix \\
Multibyte chunker dropped six characters & Repair the chunker; re-run the affected step \\
Sim.\ vs.\ sim.\ (VVUQ-02 limits) & Stage toward physical validation under standards \\
\bottomrule
\end{tabularx}
\caption{Limitations from the codegen and execution steps, each with a remedy.}
\label{tab:limits-remedies}
\end{table}

% \clearpage
